Virtual unrolling of the Herculaneum papyri produces surface meshes that sometimes drift from one winding of the scroll to the next. Such sheet switches are the dominant failure mode of automated segmentation and are usually caught by eye. This note reports an end-to-end, logged, CPU-only pipeline for PHerc1447 (seed selection, mesh growth, surface rendering, cross-section QA and ink inference on Apple Silicon) together with two geometric quality measures borrowed from physics. The first treats the mesh as a Riemannian surface and asks whether it is developable: the holonomy of parallel transport around small loops, computed as a sum of discrete angle deficits via Gauss--Bonnet, should vanish for a sheet of papyrus. It does, even across a sheet switch, and we show why: the switch is extrinsic, so intrinsic curvature cannot see it. The second treats the papyrus weave as a two-dimensional tetratic phase and measures its order parameter $\psifour=\langle e^{4i\theta}\rangle$ and the quantised disclination charge of its orientation field. Tetratic order separates a sheet switch from clean weave where the standard structure-tensor coherence does not. We also calibrate the released ink-detection models against the team's PHerc.172 w087 map and report a depth-window scan with block-bootstrap confidence intervals. Code, meshes, logs and figures are public.
